\section{Additional Exercises}

\subsection{Geodesics In A Schwarzschild Spacetime}

\begin{exercise}\label{ex:schwarzschild-1}
The Schwarzschild metric is given and we are told to use the light hand notation
\bse
    t(\lambda) := \big(x^0\circ\gamma)(\lambda),
\ese
and similarly for $r(\lambda)$, $\theta(\lambda)$ and $\varphi(\lambda)$, where $\gamma:\R\to U$ is some curve.

Write down the Lagrangian $\cL := g_{ab}\dot{\gamma}^a\dot{\gamma}^b$!
\end{exercise}
\hyperref[sol:schwarzschild-1]{\small \textit{Solution on p.~\pageref*{sol:schwarzschild-1}}}

\begin{exercise}\label{ex:schwarzschild-2}
Find the Euler-Lagrange equation with respect to $t(\lambda)$!
\end{exercise}
\hyperref[sol:schwarzschild-2]{\small \textit{Solution on p.~\pageref*{sol:schwarzschild-2}}}

\begin{exercise}\label{ex:schwarzschild-3}
Show that the Lie derivative of $g$ with respect to the vector field $K_t := \frac{\p}{\p t}$ vanishes. What does this mean?
\end{exercise}
\hyperref[sol:schwarzschild-3]{\small \textit{Solution on p.~\pageref*{sol:schwarzschild-3}}}

\begin{exercise}\label{ex:schwarzschild-4}
The exact form of the conserved quantity is given by $(K_t)_a(x^a)'(\lambda)=$const. (without proof). Derive an expression for the quantity $t'(\lambda)$ appearing in the Lagrangian!
\end{exercise}
\hyperref[sol:schwarzschild-4]{\small \textit{Solution on p.~\pageref*{sol:schwarzschild-4}}}

\begin{exercise}\label{ex:schwarzschild-5}
Moreover, we can find so-called "spherical symmetry", that is, the Lie derivative of $g$ with respect to the already known vector fields
\bse
    \begin{split}
        X_1 & = \sin\varphi\frac{\p}{\p\theta} + \cot\theta\cos\varphi\frac{\p}{\p\varphi} \\
        X_2 & = \cos\varphi\frac{\p}{\p\theta} - \cot\theta\sin\varphi\frac{\p}{\p\varphi} \\
        X_3 & = \frac{\p}{\p\varphi}
    \end{split}
\ese
vanishes. What physical quantity is conserved by this symmetry?
\end{exercise}
\hyperref[sol:schwarzschild-5]{\small \textit{Solution on p.~\pageref*{sol:schwarzschild-5}}}

\begin{exercise}\label{ex:schwarzschild-6}
Due to $X_1$ and $X_2$ (without proof), one can fix the motion to a plane of constant $\theta=\frac{\pi}{2}$. How can you derive an expression for the remaining term $\varphi'(\lambda)$?
\end{exercise}
\hyperref[sol:schwarzschild-6]{\small \textit{Solution on p.~\pageref*{sol:schwarzschild-6}}}

\begin{exercise}\label{ex:schwarzschild-7}
Use all the fact that $\cL=1$ on the parameterisation. Insert the previously obtained results and take all terms not containing $E$ to one side!
\end{exercise}
\hyperref[sol:schwarzschild-7]{\small \textit{Solution on p.~\pageref*{sol:schwarzschild-7}}}

\begin{exercise}\label{ex:schwarzschild-8}
Can you interpret the terms appearing in this expression?
\end{exercise}
\hyperref[sol:schwarzschild-8]{\small \textit{Solution on p.~\pageref*{sol:schwarzschild-8}}}

\subsection{Gravitational Redshift}

\begin{exercise}\label{ex:schwarzschild-9}
Consider a spacetime equipped with the Schwarzschild metric as well as two observers 1 and 2 at rest in their respective system of reference ($\dot{r}=0$, $\dot{\theta}=0$, $\dot{\varphi}=0$). The observers sit at the same $\theta$ and $\varphi$ while $r_1<r_2$.

Derive an expression for $t'(\lambda)$ using the Lagrangian from the previous exercise!
\end{exercise}
\hyperref[sol:schwarzschild-9]{\small \textit{Solution on p.~\pageref*{sol:schwarzschild-9}}}

\begin{exercise}\label{ex:schwarzschild-10}
Observer 1 emits photons that observer 2 detects. The gap between the two photon emissions is $\Delta \lambda_1$. Find the gap $\Delta \lambda_2$ seen by observer 2!
\end{exercise}
\hyperref[sol:schwarzschild-10]{\small \textit{Solution on p.~\pageref*{sol:schwarzschild-10}}}

\begin{exercise}\label{ex:schwarzschild-11}
Consider the ratio of frequencies $\frac{\omega_1}{\omega_2}$. What happens for observer 2 being approximately at infinity? What happens when sending $r_1$ to the Schwarzschild radius $r_s=2GM$?
\end{exercise}
\hyperref[sol:schwarzschild-11]{\small \textit{Solution on p.~\pageref*{sol:schwarzschild-11}}}
